\begin{scanpage}{0176}{158}
to an integral equation for a ``shadow'' current which is, in our terminology, half the nonuniform part of the current. The resulting integral equations may be solved (with any ratio between the dimensions of the hole and the wavelength) by the method of successive approximations. Moreover, they allow one to obtain asymptotic expressions which are suitable for short waves. In Reference [55] Greenberg found an asymptotic expression for the current on a strip with $ka\gg1$ ($2a$ is the strip's width). Greenberg and Pimenov [56] obtained a similar solution in the case of normal incidence of a plane wave on a circular hole. Using the same method, an asymptotic expression was found for the current on a flat ring [57], the width and inner diameter of which are a great deal larger than the wavelength.

The above listed works [53--57] already relate to the mathematical theory of diffraction: in them the first terms of the asymptotic expansions for the current were obtained with the desire evidently to also be able to calculate the following terms. Unfortunately, the asymptotic expressions which have been found up to now refer only to currents, and one is obliged to calculate the scattering characteristics by means of numerical quadratures [56]. As a consequence of the rapid oscillation of the integrands, such a method leads to rather unwieldy calculations and does not enable one to formulate a clear representation of the fringing field formation, and also does not allow one to study this field properly.

Millar [58] investigated the problems of electromagnetic wave diffraction by slits in a flat screen. The system of integral equations obtained by him for the current is solved by the method of successive approximations. The field in the hole is calculated from the currents which are found, and then on the basis of the field in the hole the field in the far zone and the transmission coefficient are calculated. All the indicated quantities are represented in the form of an asymptotic expansion in reciprocal powers of the parameter $\sqrt{ka}$. A solution also is obtained in the case of glancing incidence of a plane wave.
\end{scanpage}
